% vim: spl=pt
\begin{exercício}{Espectro do adjunto e do inverso de um operador limitado}{ex7}
  Seja \(T \in \bounded(\hilbert).\) Mostre que
  \begin{enumerate}[label=(\alph*)]
    \item \(\sigma(T^*) = \setc{\conj{\lambda} \in \mathbb{C}}{\lambda \in \sigma(T)};\) e
    \item se \(0 \notin \sigma(T),\) \(T^{-1} \in \bounded(\hilbert)\) existe e \(\sigma(T^{-1}) = \setc{\lambda^{-1} \in \mathbb{C}}{\lambda \in \sigma(T)}.\)
  \end{enumerate}
\end{exercício}
\begin{proof}[Resolução de (a)]
  Seja \(\mathrm{Inv}(\hilbert)\) o conjunto de operadores limitados e invertíveis agindo em \(\hilbert\), isto é, \(A \in \mathrm{Inv}(\hilbert)\) se \(A : \hilbert \to \hilbert\) é um operador limitado bijetivo. Notemos que \(\mathrm{Inv}(\hilbert)\) é um subconjunto auto-adjunto de \(\bounded(\hilbert)\), já que se \(A \in \mathrm{Inv}(\hilbert),\)  então \((A^*)^{-1} = (A^{-1})^*,\) logo \(A^* \in \mathrm{Inv}(\hilbert).\) Assim,
  \begin{equation*}
    \lambda \in \sigma(T) \iff \lambda \unity - T \notin \mathrm{Inv}(\hilbert) \iff \conj{\lambda}\unity - T^* \notin \mathrm{Inv}(\hilbert) \iff \conj{\lambda} \in \sigma(T^*),
  \end{equation*}
  logo \(\sigma(T^*) = \setc{\conj{\lambda} \in \mathbb{C}}{\lambda \in \sigma(T)}.\)
\end{proof}
\begin{proof}[Resolução de (b)]
  Notemos que \(\mathrm{Inv}(\hilbert)\) é um grupo sob composição, já que o elemento neutro é \(\unity,\) a composição é associativa, a inversa dos elementos é garantida por definição, e é fechado sob composição, já que se \(A, B \in \mathrm{Inv}(\hilbert),\) então \(AB \in \mathrm{Inv}(\hilbert)\) com \((AB)^{-1} = B^{-1} A^{-1}.\) Ainda, \(\mathrm{Inv}(\hilbert)\) é fechado sob multiplicação por \(\mathbb{C} \setminus \set{0}\), já que se \(A \in \mathrm{Inv}(\hilbert),\) então \(\alpha A \in \mathrm{Inv}(\hilbert)\) sempre que \(\alpha \neq 0,\) já que \((\alpha A)^{-1} = \frac{1}{\alpha} A^{-1}.\)

  Se \(0 \notin \sigma(T),\) então \(0 \in \rho(T),\) logo \(T\) é invertível com \(T^{-1} \in \bounded(\hilbert),\) e, por conseguinte, \(0 \notin \sigma(T^{-1}).\) Seja \(\lambda \in \rho(T) \setminus \set{0},\) então como
  \begin{equation*}
    (\lambda \unity - T) T^{-1} = -(\unity - \lambda T^{-1}),
  \end{equation*}
  concluímos que
  \begin{equation*}
    \lambda \unity - T \in \mathrm{Inv}(\hilbert) \iff \unity - \lambda T^{-1} \in \mathrm{Inv}(\hilbert).
  \end{equation*}
  De fato, \(\lambda \unity - T \in \mathrm{Inv}(\hilbert)\) e \(T^{-1} \in \mathrm{Inv}(\hilbert),\) portanto \(\unity - \lambda T^{-1}\) é um elemento de \(\mathrm{Inv}(\hilbert).\) Assim, segue que
  \begin{align*}
    \lambda \in \sigma(T) &\iff \lambda \neq 0 \land \lambda\unity - T \notin \mathrm{Inv}(\hilbert)\\
                          &\iff \lambda \neq 0 \land \unity - \lambda T^{-1} \notin \mathrm{Inv}(\hilbert)\\
                          &\iff \lambda \neq 0 \land \frac{1}{\lambda} \unity - T^{-1} \notin \mathrm{Inv}(\hilbert)\\
                          &\iff \lambda \neq 0 \land \frac{1}{\lambda} \in \sigma(T^{-1})\\
                          &\iff \frac{1}{\lambda} \in \sigma(T^{-1})
  \end{align*}
  isto é,
  \begin{equation*}
    \sigma(T^{-1}) = \setc{\lambda^{-1} \in \mathbb{C}}{\lambda \in \sigma(T)},
  \end{equation*}
  como desejado.
\end{proof}
